\section{Governed Semantic Routing and Executable Falsification}
\label{sec:semantic-routing-executable-falsification}

This cycle tested an outcome-governed semantic-memory router and executable,
model-generated falsification programs. Neither challenger was promoted. These
negative decisions distinguish useful infrastructure and partial falsification
competence from a stable retained capability.

\subsection{Semantic-memory routing}

The router compared source-and-issue-only context, documentation-derived
ontology memory, and expanded source/test/call-site memory. Unsafe or regressing
policies were excluded and the minimum-cost safe policy was selected.

\begin{table}[ht]
\centering
\caption{Governed semantic-router result.}
\label{tab:semantic-router-results}
\begin{tabular}{lr}
\hline
Measure & Result \\
\hline
Outcome rows & 9 \\
Context policies & 3 \\
Repositories & 3 \\
Routed success & 100\% \\
Weakest-repository success & 100\% \\
Routed proposal attempts & 4 \\
Never-memory attempts & 5 \\
Always-documentation attempts & 4 \\
Unsafe acceptances & 0 \\
\hline
\end{tabular}
\end{table}

The router selected documentation memory. It beat never-memory but tied the
always-documentation champion, and did not meet the five-repository diversity
gate. CAU rejected
\texttt{procedure\_governed\_semantic\_memory\_router\_4e76e58498a8}.

\subsection{Executable falsification}

For each selected repair the proposal substrate generated a standalone
functional program and an adversarial boundary/security program. No historical
human patch or hidden verifier test was included in generation. An
abstract-syntax-tree audit rejected network, subprocess, dynamic-execution,
file-write, permission-changing and verification-bypass operations. A
repository passed only if
\begin{equation}
 \neg T_f(x_{\mathrm{original}}) \land
 T_f(x_{\mathrm{candidate}}) \land
 T_a(x_{\mathrm{candidate}}),
\end{equation}
where $T_f$ is the functional falsifier and $T_a$ is the adversarial program.

\begin{table}[ht]
\centering
\caption{Executable-test invention and one outcome-criticism round.}
\label{tab:executable-falsification-results}
\begin{tabular}{lrr}
\hline
Measure & Initial & Revised \\
\hline
Original buggy repositories rejected & 3/3 & 3/3 \\
Selected repairs accepted & 1/3 & 1/3 \\
Adversarial programs accepted & 1/3 & 1/3 \\
End-to-end success & 33.33\% & 33.33\% \\
Unsafe programs executed & 0 & 0 \\
Timeouts & 0 & 0 \\
Live writes & 0 & 0 \\
\hline
\end{tabular}
\end{table}

All generated functional programs falsified their original buggy checkout.
However, pydicom did not reach the required lazy materialisation path and pvlib
failed during an unrelated package-level NumPy incompatibility. A trace-only
criticism turn diagnosed both failures but did not repair the execution
harnesses. CAU therefore rejected
\texttt{procedure\_executable\_falsification\_invention\_76dbafce2714} and
\texttt{procedure\_revised\_executable\_falsification\_e4d6d473c3db}.

\subsection{Resulting design requirement}

Property invention must be separated from execution-environment grounding. The
next system should acquire a minimal verified adapter for each unfamiliar
runtime or test framework, prove that it reaches the target behavior, and then
generate counterexamples and security properties inside that adapter.
Promotion remains gated on original-fail/candidate-pass behavior, malicious
candidate rejection, repeated-seed stability, at least five repositories and
three languages, followed by independently administered hidden evaluation.

This is partial executable falsification on a public bounded cohort, not
reliable autonomous verification or independent external certification.
